\chapter{Introduction}\label{ch:intro}

\section*{Project Overview}

Capacity planning is a core topic in performance engineering, focusing on
how systems behave under changing workloads and resource allocations.
The COMP9334 course at UNSW introduces students to queueing theory,
operational laws, and discrete-event simulation (DES), but historically
lacks practical, hands-on experimentation using real systems. Without
concrete measurements, students often struggle to build intuition about
how theoretical models manifest in real-world environments.

This project addresses that gap by developing a low-cost, reproducible
platform that allows students to run controlled experiments on
containerised services, collect performance metrics, and compare the
results directly against analytical and simulation-based predictions.

\section*{Project Aims}

The project aims to:

\begin{itemize}
    \item Provide accessible, container-based experimental infrastructure
          that students can run on their own machines.
    \item Enable hands-on exploration of queueing models, utilisation
          laws, tail-latency behaviour, and capacity limits.
    \item Allow students to collect metrics from real services and use
          them to parameterise analytical models and DES.
    \item Demonstrate how empirical traces can support predictive
          performance modelling for various resource configurations.
\end{itemize}

\section*{System Requirements and Approach}

To be effective for teaching, the platform must satisfy several
requirements:

\begin{itemize}
    \item \textbf{Ease of deployment:} The environment should run on
          Windows, macOS, or Linux with minimal setup, using Docker
          Compose to orchestrate components.
    \item \textbf{Lightweight resource usage:} Experiments should be small
          enough to run on student laptops without specialised hardware.
    \item \textbf{Configurable workloads and resources:} Students should be
          able to vary CPU limits, memory quotas, and arrival rates to
          study performance under scaling and saturation.
    \item \textbf{Integrated monitoring:} Tools such as cAdvisor,
          Prometheus, and Grafana provide real-time visibility of CPU
          usage, memory behaviour, and overall utilisation.
    \item \textbf{Modelling compatibility:} Collected traces must be usable
          as inputs for queueing analysis, operational laws, and
          trace-driven DES.
\end{itemize}

This approach creates a clear bridge between theory and practice:
students observe how latency, throughput, and utilisation behave in a
real system while applying the analytical methods they learn in class.

\bigskip

The remainder of this thesis is organised as follows.  
Chapter~\ref{ch:background} introduces the theoretical background behind
capacity planning and relevant modelling techniques.  
Chapter~\ref{ch:eval} presents the experimental results and validates the
behaviour of the platform.  
Chapter~\ref{ch:conclusion} summarises key findings and outlines the next
steps planned for Thesis~B.

